\%============================================================
\chapter{Bài Học: Bạn Trong Thù, Thù Trong Bạn (Friend Hidden in Enemy, Enemy Hidden in Friend)}
\begin{center}
  \textit{\color{truyengold}``The harshest critic we have is sometimes our greatest teacher.''}\\[4pt]
  \textit{(Người chỉ trích ta gay gắt nhất đôi khi là người thầy tốt nhất của ta.)}\\[4pt]
  \small\color{truyenblue}--- Cổ Kim Đông Tây ---
\end{center}
\ngancach

\section{Lincoln Biến Kẻ Thù Thành Nội Các}
\begin{truyen}{Lincoln Biến Kẻ Thù Thành Nội Các}{Abraham Lincoln --- Hoa Kỳ, 1860}
\chuhoa{K}{hi} Lincoln đắc cử Tổng thống năm 1860, ông có ba đối thủ tranh cử cực kỳ gay gắt: William Seward, Salmon Chase, và Edward Bates.\\[4pt]
\textit{(When Lincoln won the presidency in 1860, he had three fiercely competitive rivals: William Seward, Salmon Chase, and Edward Bates.)}

Mỗi người trong số đó đều cho rằng mình xứng đáng làm tổng thống hơn Lincoln và không che dấu điều đó.\\[4pt]
\textit{(Each of them believed he was more qualified than Lincoln and made no secret of it.)}

Thế nhưng Lincoln bổ nhiệm cả ba người vào nội các của mình: Seward làm Ngoại trưởng, Chase làm Bộ trưởng Tài chính, Bates làm Bộ trưởng Tư pháp.\\[4pt]
\textit{(Yet Lincoln appointed all three into his cabinet: Seward as Secretary of State, Chase as Treasury Secretary, Bates as Attorney General.)}

Cộng sự hỏi tại sao ông lại chọn những người ghét ông nhất vào nội các; ông trả lời: ``Vì họ là những người tài nhất đất nước --- và đất nước cần người tài, không cần người yêu mến tôi.''\\[4pt]
\textit{(An aide asked why he chose his greatest enemies for the cabinet; he answered: `Because they are the most capable men in the country --- and the country needs capable men, not men who love me.')}

Nội các ``kẻ thù'' đó lãnh đạo đất nước qua Nội chiến và giúp Lincoln để lại di sản vĩ đại nhất trong lịch sử Mỹ.\\[4pt]
\textit{(That cabinet of `enemies' led the nation through the Civil War and helped Lincoln leave the greatest legacy in American history.)}

\end{truyen}
\begin{baihoc}
  Người thù bạn thường biết điểm yếu của bạn rõ hơn bạn bè --- hãy tận dụng điều đó.\\[4pt]
  \textit{(Those who hate you often know your weaknesses better than your friends --- use that.)}
\end{baihoc}

\section{Trần Quốc Tuấn Và Anh Giết Cha}
\begin{truyen}{Trần Quốc Tuấn Và Anh Giết Cha}{Trần Hưng Đạo --- Việt Nam, thế kỷ 13}
\chuhoa{T}{rần} Quốc Tuấn lớn lên với thù nhà sâu đậm: cha ông --- Trần Liễu --- bị anh là Trần Thái Tông đoạt vợ và quyền lực, chết trong cay đắng.\\[4pt]
\textit{(Tran Quoc Tuan grew up with deep family hatred: his father Tran Lieu had his wife and power seized by his brother Tran Thai Tong, and died in bitterness.)}

Cha ông trăng trối bảo ông phải trả thù, giành lại danh dự cho dòng họ.\\[4pt]
\textit{(His dying father charged him to take revenge and restore family honor.)}

Nhưng khi quân Nguyên Mông tràn vào, Trần Quốc Tuấn hiểu rằng nếu cứ ôm hận thù trong gia tộc thì cả đất nước sẽ mất.\\[4pt]
\textit{(But when the Mongols invaded, Tran Quoc Tuan understood that clinging to family revenge would cost the entire nation.)}

Ông gác thù riêng lại, dâng hết tài năng phụng sự triều đình và chỉ huy quân dân đánh đuổi ba lần xâm lược Nguyên Mông.\\[4pt]
\textit{(He set aside personal hatred, dedicated all his talents to serving the court, and commanded the people to defeat three Mongol invasions.)}

Sau thắng lợi, vua hỏi ông về lời trăng trối của cha; ông quỳ xuống và nói: ``Thần đã dâng lời hứa với cha lên làm lễ vật tế cho đất nước.''\\[4pt]
\textit{(After victory, the king asked about his father's deathbed words; he knelt and said: `I offered my vow to my father as a sacrifice for the nation.')}

\end{truyen}
\begin{baihoc}
  Kẻ thù thật sự không phải là người ở ngoài kia --- mà là lòng thù hận ở trong ta ngăn ta làm điều vĩ đại.\\[4pt]
  \textit{(The real enemy is not the one out there --- it is the hatred within us that prevents us from doing great things.)}
\end{baihoc}

\section{Oprah Winfrey Và Người Quản Lý Tàn Nhẫn}
\begin{truyen}{Oprah Winfrey Và Người Quản Lý Tàn Nhẫn}{Oprah Winfrey --- Hoa Kỳ, 1976}
\chuhoa{K}{hi} còn trẻ, Oprah Winfrey làm phóng viên truyền hình ở Baltimore; người quản lý của cô nhận xét thẳng mặt: ``Cô không phù hợp với truyền hình.''\\[4pt]
\textit{(Early in her career, Oprah Winfrey worked as a TV reporter in Baltimore; her producer told her to her face: `You are not suitable for television.')}

Ông ta chuyển cô xuống làm dẫn chương trình buổi sáng --- mà ông nghĩ là thứ cỏn con không quan trọng.\\[4pt]
\textit{(He moved her down to hosting a morning talk show --- which he thought was a trivial, unimportant assignment.)}

Nhưng trên ghế dẫn chương trình đó, Oprah khám phá ra điểm mạnh thật sự của mình: khả năng kết nối cảm xúc với khán giả.\\[4pt]
\textit{(But in that hosting chair, Oprah discovered her true strength: the ability to connect emotionally with her audience.)}

Chương trình buổi sáng của cô nhanh chóng thống trị kênh địa phương và đưa cô thẳng đến The Oprah Winfrey Show --- chương trình truyền hình được xem nhiều nhất lịch sử Mỹ.\\[4pt]
\textit{(Her morning show quickly dominated local ratings and launched her directly to The Oprah Winfrey Show --- the most-watched talk show in American television history.)}

Người đã ``phạt'' cô bằng cách chuyển xuống ghế dẫn nhỏ --- thực ra đã trao cho cô cánh cửa cuộc đời.\\[4pt]
\textit{(The person who `punished' her by moving her to a small hosting seat --- had actually given her the door to her life.)}

\end{truyen}
\begin{baihoc}
  Người cho bạn cái mà bạn không muốn đôi khi là người cho bạn thứ bạn cần nhất.\\[4pt]
  \textit{(The person who gives you what you don't want is sometimes the one giving you what you need most.)}
\end{baihoc}

\section{Gandhi Và Người Anh Bỏ Tù Ông}
\begin{truyen}{Gandhi Và Người Anh Bỏ Tù Ông}{Mahatma Gandhi --- Nam Phi, 1906}
\chuhoa{K}{hi} viên tướng Anh Jan Smuts bỏ tù Gandhi lần đầu năm 1908 vì ông phản đối đạo luật phân biệt chủng tộc, ông ta cho rằng mình đã dập tắt được phong trào.\\[4pt]
\textit{(When British general Jan Smuts first jailed Gandhi in 1908 for opposing apartheid laws, he believed he had crushed the movement.)}

Nhưng mỗi lần Gandhi bị bắt giam, tin tức lan rộng hơn và phong trào đấu tranh mạnh hơn.\\[4pt]
\textit{(But every time Gandhi was jailed, the news spread further and the movement grew stronger.)}

Smuts nhận ra ông ta đang giúp Gandhi nổi tiếng và quyết định thả tự do --- nhưng đã quá muộn.\\[4pt]
\textit{(Smuts realized he was making Gandhi famous and decided to release him --- but it was already too late.)}

Về sau, Smuts viết về Gandhi: ``Đây là người đã thay đổi tôi; ông ta dạy tôi rằng sức mạnh không phải ở thanh kiếm.''\\[4pt]
\textit{(Later, Smuts wrote of Gandhi: `This is the man who changed me; he taught me that strength does not lie in the sword.')}

Hai người cựu thù --- người đã bỏ tù nhau --- cuối đời trao đổi thư từ như những người bạn tri kỷ.\\[4pt]
\textit{(The two former enemies --- who had jailed each other --- corresponded in their final years like old friends.)}

\end{truyen}
\begin{baihoc}
  Kẻ thù đủ mạnh để thách thức chúng ta đôi khi là người dạy ta bài học lớn nhất.\\[4pt]
  \textit{(An enemy strong enough to challenge us is sometimes the one who teaches our greatest lesson.)}
\end{baihoc}

\section{Steve Jobs Và Bill Gates}
\begin{truyen}{Steve Jobs Và Bill Gates}{Steve Jobs / Bill Gates --- Hoa Kỳ, 1980--2007}
\chuhoa{S}{teve} Jobs và Bill Gates là đối thủ thương mại khốc liệt nhất của thế kỷ 20; họ công khai chỉ trích nhau trong nhiều thập kỷ.\\[4pt]
\textit{(Steve Jobs and Bill Gates were the fiercest business rivals of the 20th century; they publicly criticized each other for decades.)}

Jobs gọi Microsoft là ``kẻ ăn cắp ý tưởng''; Gates gọi Jobs là người tạo ảo giác hơn là tạo sản phẩm thực.\\[4pt]
\textit{(Jobs called Microsoft `a thief of ideas'; Gates called Jobs a man who created illusions more than real products.)}

Nhưng ít ai biết rằng chính Bill Gates đã cứu Apple bằng khoản đầu tư 150 triệu đô năm 1997 khi Apple gần phá sản.\\[4pt]
\textit{(But few knew that Bill Gates saved Apple with a 150-million-dollar investment in 1997 when Apple was near bankruptcy.)}

Và Jobs thừa nhận rằng chính sự cạnh tranh với Microsoft đã buộc Apple phải không ngừng đổi mới.\\[4pt]
\textit{(And Jobs admitted that competition with Microsoft was what forced Apple to innovate constantly.)}

Khi Jobs hấp hối năm 2011, Gates là một trong số ít người ông nói chuyện lần cuối --- hai kẻ thù vĩ đại nhất ngành công nghệ hóa ra cũng là những người hiểu nhau nhất.\\[4pt]
\textit{(As Jobs lay dying in 2011, Gates was one of the few people he spoke to at the end --- the two greatest enemies in tech turned out to be the ones who understood each other best.)}

\end{truyen}
\begin{baihoc}
  Đối thủ thực sự là người duy nhất đủ tầm để thách thức bạn lớn lên; không có họ, bạn có thể không biết mình có thể đi xa đến đâu.\\[4pt]
  \textit{(A true rival is the only one at your level who can challenge you to grow; without them, you may never know how far you can go.)}
\end{baihoc}

\section{Tôn Tử Và Tướng Trên Chiến Trường}
\begin{truyen}{Tôn Tử Và Tướng Trên Chiến Trường}{Tôn Tử --- Trung Quốc, thế kỷ 5 TCN}
\chuhoa{T}{ôn} Tử trong binh thư viết: ``Hãy biết địch như biết bạn --- vì địch cho bạn biết điều bạn thiếu.''\\[4pt]
\textit{(Sun Tzu wrote in the Art of War: `Know your enemy as your friend --- for the enemy shows you what you lack.')}

Ông kể chuyện một vị tướng luôn thắng trận nhờ thói quen học từ địch: sau mỗi trận, ông phân tích kỹ những đòn tấn công mà địch dùng hiệu quả.\\[4pt]
\textit{(He told of a general who always won battles by learning from the enemy: after each battle, he carefully analyzed whatever attacks the enemy used effectively.)}

Vị tướng đó không coi kẻ địch như kẻ thù mà như người thầy không mời --- người tự nguyện chỉ ra điểm yếu.\\[4pt]
\textit{(That general did not see the enemy as a foe but as an uninvited teacher --- one who voluntarily pointed out his weaknesses.)}

Quân đội của ông sau nhiều năm chiến đấu trở thành đội quân biết chiến thuật đa dạng nhất --- vì họ học từ nhiều địch thủ nhất.\\[4pt]
\textit{(His army, after years of fighting, became the force with the most varied tactics --- because they had learned from the most diverse enemies.)}

Bài học Tôn Tử: kẻ địch giỏi nhất là người thầy giỏi nhất --- nếu bạn đủ khiêm nhường để học.\\[4pt]
\textit{(Sun Tzu's lesson: the best enemy is the best teacher --- if you are humble enough to learn.)}

\end{truyen}
\begin{baihoc}
  Thay vì tiêu diệt kẻ thù, hãy học được từ họ trước --- điều đó vừa khôn ngoan hơn vừa có ích hơn.\\[4pt]
  \textit{(Before destroying your enemy, learn from them first --- that is both wiser and more profitable.)}
\end{baihoc}

\section{Nguyễn Huệ Học Binh Pháp Từ Quân Xiêm Thua Trận}
\begin{truyen}{Nguyễn Huệ Học Binh Pháp Từ Quân Xiêm Thua Trận}{Quang Trung --- Việt Nam, 1785}
\chuhoa{S}{au} khi đánh tan quân Xiêm tại Rạch Gầm -- Xoài Mút năm 1785, Nguyễn Huệ không chỉ mừng chiến thắng.\\[4pt]
\textit{(After destroying the Siamese army at Rach Gam -- Xoai Mut in 1785, Nguyen Hue did not simply celebrate the victory.)}

Ông triệu tập tướng sĩ nghiên cứu kỹ các chiến thuật mà quân Xiêm đã dùng và những điểm mà họ gần như đã thắng.\\[4pt]
\textit{(He summoned his officers to carefully study the tactics the Siamese had used and the moments when they had nearly won.)}

Ông nói: ``Biết ta thắng không đủ; phải biết tại sao địch gần thắng.''\\[4pt]
\textit{(He said: `Knowing we won is not enough; we must know why the enemy nearly won.')}

Những bài học từ quân Xiêm được tích hợp vào chiến lược của quân Tây Sơn --- chuẩn bị cho trận Đống Đa lẫy lừng năm 1789.\\[4pt]
\textit{(The lessons from the Siamese forces were incorporated into the Tay Son army's strategy --- preparing for the brilliant Battle of Dong Da in 1789.)}

Tại Đống Đa, Quang Trung đại phá 29 vạn quân Thanh chỉ trong năm ngày đêm --- chiến thắng thần tốc bậc nhất lịch sử châu Á.\\[4pt]
\textit{(At Dong Da, Quang Trung destroyed 290,000 Qing troops in just five days and nights --- the fastest decisive victory in Asian history.)}

\end{truyen}
\begin{baihoc}
  Chiến thắng của hôm nay nên là bài học cho ngày mai --- và bài học hay nhất đến từ chính kẻ địch.\\[4pt]
  \textit{(Today's victory should be tomorrow's lesson --- and the best lessons come from the enemy.)}
\end{baihoc}

\section{Churchill Và Hitler --- Kẻ Thù Tạo Ra Lãnh Tụ}
\begin{truyen}{Churchill Và Hitler --- Kẻ Thù Tạo Ra Lãnh Tụ}{Winston Churchill --- Anh, 1940}
\chuhoa{W}{inston} Churchill bị coi là chính khách lỗi thời và thất bại trong suốt những năm 1930; ông không được ai nghe.\\[4pt]
\textit{(Winston Churchill was considered an outdated and failed politician throughout the 1930s; no one listened to him.)}

Nhưng khi Hitler bắt đầu xâm lược châu Âu và các chính khách khác muốn nhượng bộ, Churchill là người duy nhất đã nói đúng.\\[4pt]
\textit{(But when Hitler began invading Europe and other politicians wanted to appease him, Churchill was the only one who had been right all along.)}

Chính mối đe dọa của Hitler đã đặt Churchill vào đúng vị trí lịch sử mà ông sinh ra để làm: Thủ tướng thời chiến của nước Anh.\\[4pt]
\textit{(Hitler's very threat placed Churchill in exactly the historical role he was born to fill: Britain's wartime Prime Minister.)}

Churchill nói: ``Nếu không có Hitler, chắc tôi đã chết già trong sự lãng quên.''\\[4pt]
\textit{(Churchill said: `Without Hitler, I would probably have died in obscurity.')}

Kẻ thù lớn nhất của Churchill đã trao cho ông sân khấu lớn nhất của cuộc đời ông --- và qua đó, Churchill trao cho thế giới can đảm để tiếp tục chiến đấu.\\[4pt]
\textit{(Churchill's greatest enemy gave him the greatest stage of his life --- and through it, Churchill gave the world the courage to keep fighting.)}

\end{truyen}
\begin{baihoc}
  Đôi khi kẻ thù lớn nhất là người duy nhất đủ lớn để gọi ra kẻ vĩ đại trong bạn.\\[4pt]
  \textit{(Sometimes the greatest enemy is the only one large enough to call forth the greatness within you.)}
\end{baihoc}

\section{Nguyễn Công Trứ Bị Giáng Chức Nhiều Lần}
\begin{truyen}{Nguyễn Công Trứ Bị Giáng Chức Nhiều Lần}{Nguyễn Công Trứ --- Việt Nam, thế kỷ 19}
\chuhoa{N}{guyễn} Công Trứ trong cuộc đời làm quan bị giáng chức, cách chức, và phát lưu nhiều lần vì dám nói thẳng sự thật.\\[4pt]
\textit{(Nguyen Cong Tru was demoted, dismissed, and exiled many times during his official career for daring to speak truth directly.)}

Những kẻ đẩy ông xuống --- dù là vua hay quyền thần --- thực ra không thể nào dập tắt được tiếng thơ và tinh thần của ông.\\[4pt]
\textit{(Those who pushed him down --- whether king or powerful ministers --- could never actually silence his poetry and spirit.)}

Mỗi lần bị đẩy khỏi triều đình, ông về quê khai hoang, trồng lúa, viết thơ --- và học thêm về dân tình.\\[4pt]
\textit{(Every time he was expelled from court, he returned home to clear land, grow rice, write poetry --- and learn more about the people's lives.)}

Chính những kẻ thù chính trị đã vô tình biến ông từ một vị quan thành người gần dân nhất, hiểu đời nhất.\\[4pt]
\textit{(Those political enemies inadvertently transformed him from an official into the one closest to the people, most understanding of life.)}

Di sản thơ Nguyễn Công Trứ sống mãi đến hôm nay --- và chúng được sinh ra không phải từ những tháng ngày quyền cao chức trọng, mà từ những năm bị đẩy xuống đáy.\\[4pt]
\textit{(Nguyen Cong Tru's poetic legacy lives to this day --- and it was born not from his years of high rank, but from his years of being pushed to the bottom.)}

\end{truyen}
\begin{baihoc}
  Kẻ thù hạ bệ bạn đôi khi là người xây nên di sản bất hủ cho bạn.\\[4pt]
  \textit{(Enemies who topple you sometimes build your most lasting legacy.)}
\end{baihoc}

\section{Phật Và Ma Vương Māra}
\begin{truyen}{Phật Và Ma Vương Māra}{Siddhattha Gotama --- Ấn Độ, thế kỷ 5 TCN}
\chuhoa{T}{heo} truyền thống Phật giáo, khi Siddhattha Gotama ngồi thiền dưới cội Bồ Đề, Ma Vương Māra kéo đến với toàn bộ đội quân của mình để phá rối.\\[4pt]
\textit{(According to Buddhist tradition, as Siddhattha Gotama meditated under the Bodhi tree, the demon Mara arrived with his entire army to disturb him.)}

Māra tung ra mọi thứ: cám dỗ của dục vọng, bóng đen của sợ hãi, những mũi tên của nghi ngờ.\\[4pt]
\textit{(Mara unleashed everything: the temptations of desire, the shadows of fear, the arrows of doubt.)}

Nhưng Siddhattha không chạy trốn --- ông ngồi yên, nhìn thẳng vào Māra và nói: ``Ta thấy ngươi.''\\[4pt]
\textit{(But Siddhattha did not flee --- he sat still, looked directly at Mara and said: `I see you.')}

Khoảnh khắc được nhìn thấy, được nhận diện --- Ma Vương mất sức mạnh và tan biến.\\[4pt]
\textit{(The moment of being seen, being recognized --- the demon lost its power and dissolved.)}

Các thiền sư giải thích: Māra không phải kẻ thù từ bên ngoài --- Māra chính là những nỗi sợ, nghi ngờ, ham muốn bên trong ta; và cách duy nhất để vượt qua là nhìn thẳng vào chúng.\\[4pt]
\textit{(The masters explain: Mara is not an enemy from outside --- Mara is the fear, doubt, and desire within us; and the only way through them is to look at them directly.)}

\end{truyen}
\begin{baihoc}
  Kẻ thù thực sự của bạn thường không ở bên ngoài --- hãy tìm bên trong và nhìn thẳng vào đó.\\[4pt]
  \textit{(Your true enemy is usually not outside --- look within and face it directly.)}
\end{baihoc}
